\section{Polynomial IOPs and the IOP-to-SNARK Paradigm}

\subsection{Interactive proofs, IOPs, and polynomial IOPs}

To use polynomial commitments inside proof systems, recall the following hierarchy:
\begin{itemize}[leftmargin=2em]
    \item an \textbf{interactive proof (IP)} exchanges a small number of messages;
    \item an \textbf{interactive oracle proof (IOP)} allows the prover to send long oracles that the verifier only queries in a few places;
    \item a \textbf{polynomial IOP} is an IOP in which the prover's main oracle messages are polynomials.
\end{itemize}

A polynomial IOP for circuit satisfiability has the form
\[
\Setup(C)\to (\pp,\vp),
\]
then the prover sends oracle access to a few polynomials $f_1,\dots,f_t$, the verifier samples random challenges $r_1,\dots,r_s$, and finally checks algebraic identities involving evaluations of those polynomials.

The key point is that the verifier never needs to read entire witness polynomials. It suffices to query them at a few random points.

\subsection{From a polynomial IOP to a SNARK}

Suppose a polynomial IOP commits to $t$ polynomials and the verifier needs $q$ evaluation proofs. If we instantiate the commitment layer with a PCS, then:
\begin{itemize}[leftmargin=2em]
    \item the prover sends $t$ polynomial commitments;
    \item each queried evaluation is accompanied by a PCS opening proof;
    \item the IOP's verifier logic is run on the opened values.
\end{itemize}

If the IOP is public coin, the Fiat-Shamir transform removes interaction by deriving verifier challenges from a hash of the transcript \cite{FiatShamir1986}:
\[
 r_1 = H(C,x,\mathsf{com}_1,\dots),\qquad
 r_2 = H(C,x,\mathsf{com}_1,\dots,r_1,\dots),\quad \text{etc.}
\]

This yields a non-interactive argument in the random oracle model. Non-interactive zero knowledge in the common-reference-string model originates with Blum, Feldman, and Micali \cite{BlumFeldmanMicali1988}; Fiat-Shamir is a different route, usually modeled through a random oracle.

\subsection{Complexity bookkeeping}

If the PCS commitment size is $\mathsf{pcsCom}$ and the opening size is $\mathsf{pcsOpen}$, then the resulting proof length is roughly
\[
 t\cdot \mathsf{pcsCom} + q\cdot \mathsf{pcsOpen}.
\]
The verifier's work is roughly
\[
 q\cdot \mathrm{time}(\mathsf{PCS\ verify}) + \mathrm{time}(\text{IOP verifier}),
\]
and the prover's work is roughly
\[
 t\cdot \mathrm{time}(\mathsf{Commit}) + q\cdot \mathrm{time}(\mathsf{Open}) + \mathrm{time}(\text{IOP prover}).
\]

This formula explains why the same Plonk IOP can lead to very different systems:
\[
\begin{array}{c|c|c}
\text{PCS} & \text{Resulting flavor} & \text{Typical tradeoff}\\
\hline
\text{KZG \cite{KateZaveruchaGoldberg2010}} & \text{Aztec/JellyFish-style Plonk} & \text{short proofs, trusted setup}\\
\text{IPA/Bulletproofs \cite{Bulletproofs2018}} & \text{Halo2-style Plonkish systems} & \text{transparent, slower verifier}\\
\text{FRI \cite{BenSassonBentovHoreshRiabzev2018FRI}} & \text{Plonky2/STARK-like systems} & \text{transparent, hash-based, larger proofs}\\
\end{array}
\]


